\uuid{IvA3}
\exo7id{1887}
\titre{exo7 1887}
\auteur{maillot}
\organisation{exo7}
\datecreate{2001-09-01}
\isIndication{false}
\isCorrection{false}
\chapitre{Topologie}
\sousChapitre{Ouvert, fermé, intérieur, adhérence}
\cpge{mpsi-mp2i.nombres-reels-suites.b, pcsi.nombres-reels-suites.a, ptsi.nombres-reels-suites.a, mpsi-mp2i.nombres-reels-suites.a, pcsi.raisonnement-vocabulaire-ensembliste.c, ptsi.raisonnement-vocabulaire-ensembliste.c, bcpst1.nombres-reels, mp-mpi.topologie-espaces-normes.d, pc.espaces-vectoriels-normes.d, psi.espaces-vectoriels-normes.d, pt.fonctions-deux-trois-variables.a, mp-mpi.topologie-espaces-normes.a, pc.espaces-vectoriels-normes.a, psi.espaces-vectoriels-normes.a}
\module{Analyse}
\niveau{L2}
\difficulte{}

\contenu{
\begin{enumerate}
    \item \question{Quelles sont toutes les normes sur le $\R$-espace vectoriel $\R${} ?

On se place désormais dans $(\R,|\, .\, |)$.}
    \item \question{Quelles sont les boules ouvertes ? fermées ?}
    \item \question{Ouverts et fermés de $\R${} :
}
\begin{enumerate}
    \item \question{soit $(I_a)_{a\in A}$ une famille d'intervalles ouverts non vides de
$\R$, deux à deux disjoints. Montrer que A est au plus dénombrable.}
    \item \question{soit $O$ un ouvert de $\R$, et $a\in O$. On pose $A=\{x\in\R\ | \  x\not\in
O\text{ et } x>a\}$ et $B=\{x\in\R\ | \  x\not\in O \text{ et }
x<a\}$. Etudier l'existence de $\inf A$ et $\sup B$.}
    \item \question{en déduire que :
}
\begin{itemize}
    \item \question{tout ouvert de $\R${} est réunion d'une famille au plus dénombrable
d'intervalles ouverts}
    \item \question{tout fermé de $\R${} est réunion d'une famille au plus dénombrable
d'intervalles fermés.
}
\end{itemize}
\end{enumerate}
\end{enumerate}
}
